\documentclass[12pt,a4paper,notitlepage]{report}
\usepackage{polski}
\usepackage[utf8]{inputenc}
\usepackage[polish]{babel}
\usepackage{url}
\usepackage[pdftex]{graphicx}
\usepackage{sidecap}
\usepackage{verbatim}
\usepackage{array}
\usepackage{amsmath}
\usepackage{enumitem}
\usepackage{fancyhdr}
\usepackage[margin=1.8cm]{geometry}
\usepackage{float}
\usepackage{listings}
\usepackage{color}
\usepackage{wrapfig,booktabs}
\usepackage{caption}
\usepackage{wrapfig}

\newcommand\blfootnote[1]{%
  \begingroup
  \renewcommand\thefootnote{}\footnote{#1}%
  \addtocounter{footnote}{-1}%
  \endgroup
}
%remove section whitespace
\usepackage{titlesec}
\titlespacing\section{0pt}{12pt plus 4pt minus 2pt}{0pt plus 2pt minus 2pt}

\makeatletter
\newcommand*{\toccontents}{\@starttoc{toc}}
\makeatother

\renewcommand\thesection{\arabic{section}}
\newcommand{\HRule}{\rule{\linewidth}{0.5mm}} %Do tytułu

\usepackage{etoolbox} %do bibliografi
\patchcmd{\thebibliography}{\chapter*}{\section*}{}{} %żeby bibliografię nie wrzucało na nową stronę


%Do kodu C++
\usepackage{xcolor} % for setting colors

% set the default code style
\lstset{
    frame=tb, % draw a frame at the top and bottom of the code block
    tabsize=4, % tab space width
    showstringspaces=false, % don't mark spaces in strings
    numbers=left, % display line numbers on the left
    breaklines=true,                % sets automatic line breaking
    commentstyle=\color{green}, % comment color
    keywordstyle=\color{blue}, % keyword color
    stringstyle=\color{red} % string color
}
%End kod

\usepackage{amsmath} % Required for \varPsi below
\usepackage{tikz}
\usepackage{varwidth}
\usetikzlibrary{shapes,arrows}
\usepackage{pgfplots}

\usepackage{xfrac} %do ułamków np.: 1\2

\begin{document}
\renewcommand{\lstlistingname}{Kod} %Zmiana listing na Kod
\renewcommand\bibname{Literatura} %Zmiana "Bibliografia" na "Literatura"
\renewcommand{\tablename}{Tabela} %Zmiana "Tablica" na "Tabela"
%\renewcommand{\figurename}{Zrzut ekranu} %Zmiana figure na Zrzut ekranu

\begin{titlepage}
\begin{center}

\textsc{\LARGE Politechnika Poznańska}\\[0.8cm]
\textsc{\Large Wydział Elektryczny \\ Informatyka}\\[1cm]

\includegraphics[scale=1]{logopoli}
\vspace{.5cm}\\
\textsc{\large Teoria Informacji i Kodowanie}\\[0.2cm]
\textsc{\normalsize Dokumentacja Projektu}\\[.5cm]


% Author and supervisor
\noindent
\begin{minipage}[t]{0.4\textwidth}
\begin{flushleft} \large
\emph{Autorzy:}\\
Michał \textsc{Majka} \\
Nr albumu: 112679\\
Piotr \textsc{Parysek} \\
Nr albumu: 106100
\end{flushleft}
\end{minipage}%
\begin{minipage}[t]{0.4\textwidth}
\begin{flushright} \large
\emph{Prowadzący:} \\
dr inż. Ewa \textsc{Idzikowska}
\end{flushright}
\end{minipage}

%\bigskip
\vspace{2cm}
% Bottom of the page
{\large \today}
%\global\let\newpagegood\newpage %żeby nie było page break po titilepage
%\global\let\newpage\relax %żeby nie było page break po titilepage

\end{center}
\end{titlepage}

\toccontents

\section{Wstęp}
Zadaniem projektowym była implementacja algorytmu kompresji bezstratnej \newline \emph{\textbf{Run-Length Encoding}} (\textbf{RLE}).

Zadanie zrealizowaliśmy w środowisku programistycznym Qt Creator 5.5.1\cite{qt} korzystając z kompilatora GCC 4.9.1\cite{gcc}.\\
Do kontroli wersji oraz plików źródłowych wykorzystaliśmy oprogramowanie Git\cite{git}, projekt hostowaliśmy w repozytorium GitHub\cite{github}.\\
Dokumentacja została wykonana w \LaTeX \cite{latex} w programie Texmaker 4.5 \cite{program} oraz edytora online: ShareLaTeX\cite{sharelatex}. 

\section{Algorytm}
\subsection{Historia}
\textbf{Run-Length Encodings}, również znane jako \textbf{Golomb Codings}, swoje ,,podwaliny'' powstania wiąże z pracami XVII-wiecznego francuskiego matematyka \textit{Blaise'a Pascal'a} związanych z probabilistyką\cite{book}. Koncepcja kodowania powtarzających się znaków była używana od początków istnienia teorii informacji (Shannon 1949, Laemmel 1951), jednakże metodę oraz sposób kodowania wynalazł i opracował \textit{Solomon Wolf Golomb}\cite{golomb}.

\subsection{Zasada działania}
Algorytm jest relatywnie prosty $\rightarrow$ przedstawia powtarzające się wartości jako dany znak i licznik powtórzeń. Na przykład ciąg znaków:
\[pppppppuuuuuttttt......ppppoooooozzzzznnnnannnnn....pppppplllll\]
Zostaje przedstawiony w ciąg postaci:
\[p7u5t5.6p4o6z5n4an5.4p6l5\]

\begin{table}[H]
\begin{tabular}{l l}
Ciąg znaków & Liczba \\\hline
$pppppppuuuuuttttt......ppppoooooozzzzznnnnannnnn....pppppplllll$& 64\\
$p7u5t5.6p4o6z5n4an5.4p6l5$& 26 \\ \hline
& $\sfrac{26}{64}\approx 0.41$
\end{tabular}
\caption{Przykładowa kompresja znakowa}
\end{table}

\section{Opis implementacji}
Program wykonano w języku C++ z metodami dostarczonymi z Qt.
\subsection{Kodowanie}
% Define block styles
\tikzstyle{decision} = [diamond, draw, fill=green!20, text width=4.5em, text badly centered, node distance=3cm, inner sep=0pt]
\tikzstyle{block} = [rectangle, draw, fill=blue!20, text centered, rounded corners, minimum height=2.5em, execute at begin node={\begin{varwidth}{15em}}, execute at end node={\end{varwidth}}]
\tikzstyle{line} = [draw, -latex']
\subsubsection{Ustalenie znaku kodowania}    
\begin{figure}[H]
\centering
\begin{tikzpicture}[node distance = 2cm, auto]
    % Place nodes
    \node [block] (init) {Zapisanie do QMap wszystkich możliwych wartości bajtu [256 możliwości], wraz z liczbą wystąpień = 0};
    \node [block, below of=init, node distance=2.5cm] (identify) {Wylicznie wystąpień każdego bajtu i zmiana poszczególnych wartości w QMap};
    \node [block, below of=identify] (evaluate) {Wybierz wartość ,,znaku''};
    \node [block, left of=evaluate, node distance=5cm] (update) {Zmień ,,znak''};
    \node [decision, below of=update] (range) {Czy osiągniętko koniec zakresu?};
    \node [decision, below of=evaluate] (decide) {Czy liczba wystąpień znaku jest równa 0?};
    \node [block, below of=decide, node distance=3cm] (stop) {Zapisz znak w zmiennej SIGN};
    \node [block, below of=range, node distance=4.5cm] (qset) {Stwórz QSet wszyskich wystąpień dwubajtowych w danym pliku};
    \node [block, below of=qset] (choose) {Wybierz wartość dwu Bajtową ,,znaku''};
    \node [decision, below of=choose] (signdecide) {Czy ,,znak'' jest w QSet?};
    \node [block, below of=signdecide, node distance=3cm] (over) {Zapisz znak w zmiennej SIGN16};
    \node [block, right of=signdecide, node distance=5cm] (notyet) {Zmień wartość ,,znaku''};
    \node [block, below of=stop] (end) {Koniec.};
    \node [block, right of=over, node distance=5cm] (theend) {Koniec};
    % Draw edges
    \path [line] (init) -- (identify);
    \path [line] (identify) -- (evaluate);
    \path [line] (evaluate) -- (decide);
    \path [line] (decide) -- node [near start] {Nie} (range);
    \path [line] (range) -- node [near start] {Nie} (update);
    \path [line] (update) |- (evaluate);
    \path [line] (decide) -- node {Tak} (stop);
    \path [line] (range) -- node [near start] {Tak} (qset);
    \path [line] (qset) -- (choose);
    \path [line] (choose) -- (signdecide);
    \path [line] (signdecide) -- node [near start] {Nie} (over);
    \path [line] (signdecide) -- node [near start] {Tak} (notyet);
    \path [line] (notyet) |- (choose);
    \path [line] (stop) -- (end);
    \path [line] (over) -- (theend);
\end{tikzpicture}
\caption{Schemat blokowy wyszukiwania ,,znaku''.}
\end{figure}
Do wyszukania ,,znaków'' wykorzystano kontener \lstinline;QMap<quint, int>;, gdzie zmienna \lstinline;quint8; (\lstinline;unsigned byte;) wskazuje na poszczególne możliwe wartości bajta, a zmienna \lstinline;int; wskazuje ilość wystąpień danego bajta w opracowywanym pliku.
\begin{lstlisting}[language=C++, caption={Deklaracja wraz z inicjalizacją kontenera QMap<quint8, int>}]
QMap<quint8, int> IIMap;
for (quint8 i = 0; i < 255; i++) {
	IIMap.insert(i, 0);
}
\end{lstlisting}

W przypadku zdarzenia, że w pliku występują wszystkie możliwe kombinacje bajta, zamiast jedno bajtowego ,,znaku'' zostaje wprowadzony znak dwu bajtowy:
\begin{lstlisting}[language=C++, caption={Deklaracja dwu bajtowej zmiennej znakowej}]
QPair<quint8, quint8> SIGN16;
\end{lstlisting}

\subsubsection{Kodowanie}
\begin{figure}[H]
\centering
\begin{tikzpicture}[node distance = 2cm,auto]
    % Place nodes
    \node [block] (start) {Dopisz do tablicy pliku wyjściowego rozmiar ,,znaku'' oraz Bajt/y ,,znak/u''};
    \node [block, below of=start] (jeden) {Wczytaj kolejny bajt};
    \node [decision, below of=jeden] (decyzja) {Czy bajt różni się od poprzedniego?};
    \node [block, right of=decyzja, node distance=5cm] (dwa) {Dopisz Bajt\\ do Listy Elementów};
    \node [block, left of=decyzja, node distance=5cm] (trzy) {Zinkrementuj liczbę\\ wystąpień ostatniego\\ elementu listy};
    \node [block, below of=decyzja, node distance=3.5cm] (cztery) {Zapisz Listę Elementów do Tablicy bajtów wyjściowych};
    \node [block, below of=cztery] (piec) {Zapisz tablicę do pliku};
	\node [block, below of=piec] (szesc) {Koniec.};    
    
    % Draw edges
    \path [line] (start) -- (jeden);
    \path [line] (jeden) -- (decyzja);
    \path [line] (decyzja) -- node [near start] {Tak} (dwa);
    \path [line] (decyzja) -- node [near start] {Nie} (trzy);
    \path [line] (decyzja) -- node [near start] {Koniec tablicy} (cztery);
    \path [line] (dwa) |- (jeden);
    \path [line] (trzy) |- (jeden);
    \path [line] (cztery) -- (piec);
    \path [line] (piec) -- (szesc);        
\end{tikzpicture}
\caption{Schemat blokowy kodowania pliku.}
\end{figure}

Do sprawnego wczytania, zliczenia i zakodowania poszczególnych bajtów stworzono kontener \lstinline;QList; struktury \lstinline;Element;. Struktura \lstinline;Element; posiada dwa pola: \lstinline;quint8 item; $\rightarrow$ oznaczające dany bajt oraz \lstinline;quint32 value; $\rightarrow$ oznaczające ilość wystąpień (Założono, że dany bajt nie powtórzy się więcej jak $4 294 967 295$ razy).
\begin{lstlisting}[language=C++, caption={Główne struktury danych kodowania}]
struct Element {
	quint8 item;
	quint32 value;
};
QList<Element> Elements;
quint8 CurrentByte;
\end{lstlisting}

Zapisanie danych do pliku odbywa się z pośrednictwem tablicy bajtów \lstinline;QByteArray;. Analiza zapisanych znaków odbywa się poprzez przejście przez wcześniej wspomnianą listę: \lstinline;QList<Element>; i odpowiednią interpretację  wartości wystąpień danego znaku. 

Jeżeli wartość powtórzenia danego znaku nie jest większa niż minimalna długość jego zastąpienia, który ma postać ,,znak'' inicjujący, bajt powtarzany, ilość powtórzeń, to wpisywana jest ,,pierwotna postać''.

W przypadku, gdy wartość ilości powtórzeń bajtu jest większa niż maksymalna wartość jaką może osiągnąć bajt - 255 - to postać zastąpienia przybiera postać: czterech bajtów ,,znaku'', bajt powtarzany i cztery bajty licznika.
\begin{lstlisting}[language=C++, caption={Zapisanie i zakodowania danych skompresowanych do pliku}]
if (e.value < 256) {
	OutByteArray.append(SIGN);
	OutByteArray.append(e.item);
	OutByteArray.append(e.value);
} else {
	OutByteArray.append(SIGN);
    OutByteArray.append(SIGN);
    OutByteArray.append(SIGN);
    OutByteArray.append(SIGN);
    OutByteArray.append(e.item);
    QByteArray TempArray;
    TempArray = RLE::IntToHex(e.value);
    OutByteArray.append(TempArray);
}
\end{lstlisting}

\subsection{Dekodowanie}
\begin{figure}[H]
\centering
\begin{tikzpicture}[node distance = 2cm,auto]
    % Place nodes
    \node [block] (start) {Wczytaj rozmiar ,,znaku''};
    \node [block, below of=start] (zero) {Wczytaj ,,znak''};
    \node [block, below of=zero] (jeden) {Wczytaj bajt};
    \node [decision, below of=jeden] (decyzja) {Czy bajt jest rónwy ,,znakowi''?};
    \node [block, right of=decyzja, node distance=5cm] (dwa) {Dopisz Bajt do\\ tablicy bajtów wyjściowych};
    \node [block, left of=decyzja, node distance=5cm] (trzy) {Wczytaj kolejny bajt\\ -powtarzany bajt-\\ i następny bajt -licznik\\ powtórzenia. Dopisz do\\ tablicy bajtów wyjściowych\\ powtarzany bajt w ilości\\ licznika powtórzenia.};
    \node [block, below of=decyzja, node distance=3cm] (cztery) {Zapisz tablicę do pliku};
    \node [block, below of=cztery] (piec) {Koniec.};
    
    % Draw edges
    \path [line] (start) -- (zero);
    \path [line] (zero) -- (jeden);
    \path [line] (jeden) -- (decyzja);
    \path [line] (decyzja) -- node [near start] {Nie} (dwa);
    \path [line] (decyzja) -- node [near start] {Tak} (trzy);
    \path [line] (decyzja) -- node [near start] {Koniec tablicy} (cztery);
    \path [line] (dwa) |- (jeden);
    \path [line] (trzy) |- (jeden);
    \path [line] (cztery) -- (piec);      
\end{tikzpicture}
\caption{Schemat blokowy dekodowania pliku.}
\end{figure}
Dekodowanie odbywa się z pomocą analogicznych struktur / zasad / funkcji co zostały użyte podczas procesu kodowania.

\section{Użytkowanie programu}
%----- Zrobić zrzuty ekranu w domku -----
\section{Testy}
%----- Wykonać testy w domku -----
\section{Porównanie kompresji}
%----- Zrobić jakieś porównanie kompresji w domku -----

\addcontentsline{toc}{section}{Literatura}  %dodanie znacznika "Literatura" do spisu treści
\begin{thebibliography}{1}
	\bibitem{latex} Kurs \LaTeX w $\pi^e$ minut \url{http://www.fuw.edu.pl/~kostecki/kurs_latexa.pdf}.
	\bibitem{program} Program Texmaker 4.5 \url{http://www.xm1math.net/texmaker/}.
	\bibitem{sharelatex} ShareLaTeX is an online LaTeX editor \url{https://www.sharelatex.com/}.	
	\bibitem{qt} Qt \url{http://www.qt.io/}.
	\bibitem{gcc} GCC, the GNU Compiler Collection \url{https://gcc.gnu.org/}.
	\bibitem{git} Git \url{http://git-scm.com/}.
	\bibitem{github} GitHub \url{https://github.com/}.
	\bibitem{golomb} Run-length encodings - S. W. Golomb (1966); IEEE Trans Info Theory 12(3):399 \url{http://urchin.earth.li/~twic/Golombs_Original_Paper/}.
	\bibitem{book} Variable-length codes for data compression / David Salomon, London : Springer, 2007.
	\bibitem{ikony} Ikony \url{http://www.flaticon.com/}.
\end{thebibliography}

	\blfootnote{Dokument wykonany w \LaTeX \cite{latex} \\ w programie Texmaker 4.5 \cite{program}}	
\end{document}
%\begin{lstlisting}[language=C++, caption={}]
%\end{lstlisting}